\chapter*{Notes}
\addcontentsline{toc}{chapter}{Notes}
\markboth{NOTES}{NOTES}

No marker interrupts the chapters. These notes are keyed by phrase and grouped by chapter, for the reader who wants a source or a small precision the prose had no room for.

\section*{Chapter 1: The Question}

\begin{itemize}
\setlength{\itemsep}{0.5em}
\item ``The man who has fed the chicken every day throughout its life at last wrings its neck instead'' is Bertrand Russell's own image, from \textit{The Problems of Philosophy} (1912), chapter 6, ``On Induction.''
\item ``David Hume named it in 1748'' points to \textit{An Enquiry Concerning Human Understanding}; the same argument appears nine years earlier in \textit{A Treatise of Human Nature} (1739).
\end{itemize}

\section*{Chapter 2: The Library}

\begin{itemize}
\setlength{\itemsep}{0.5em}
\item ``In 1941, Jorge Luis Borges described a library'' dates the Spanish original, \textit{La biblioteca de Babel}; English readers most often meet it only in James E. Irby's 1962 translation.
\end{itemize}

\section*{Chapter 3: The Indexing}

\begin{itemize}
\setlength{\itemsep}{0.5em}
\item ``In 1946, the physicist R.\,T.\ Cox proved'' cites Cox, ``Probability, Frequency, and Reasonable Expectation,'' \textit{American Journal of Physics} 14 (1946). The standard modern development is E.\,T.\ Jaynes, \textit{Probability Theory: The Logic of Science} (2003), published five years after Jaynes's death. Appendix~\ref{app:bedrock} states the theorem precisely.
\end{itemize}

\section*{Chapter 4: The Prior}

\begin{itemize}
\setlength{\itemsep}{0.5em}
\item ``This is Kraft's inequality, proved in 1949'' cites L.\,G.\ Kraft's 1949 MIT thesis; it is sometimes called the Kraft-McMillan inequality, after Brockway McMillan's 1956 extension to all uniquely decodable codes.
\item The epigraph is Ray Solomonoff, ``A Formal Theory of Inductive Inference,'' \textit{Information and Control} 7 (1964). ``Self-delimiting techniques developed by Chaitin and Levin in the 1970s'' refers to Gregory Chaitin's 1975 \textit{Journal of the ACM} paper and Leonid Levin's independent construction of prefix complexity; Appendix~\ref{app:bedrock} states $M$ and $K$ formally.
\end{itemize}

\section*{Chapter 5: The Optimality}

\begin{itemize}
\setlength{\itemsep}{0.5em}
\item The epigraph and ``in 2005, Marcus Hutter defined AIXI'' cite Marcus Hutter, \textit{Universal Artificial Intelligence} (2005). The invariance theorem is credited jointly to Solomonoff, Kolmogorov, and Chaitin, each arriving at a version of it independently in the mid-1960s.
\end{itemize}

\section*{Chapter 6: The Slip}

\begin{itemize}
\setlength{\itemsep}{0.5em}
\item The epigraph, ``It from bit,'' is John Archibald Wheeler, ``Information, Physics, Quantum'' (1990), reprinted in \textit{Complexity, Entropy, and the Physics of Information}, edited by Wojciech Zurek.
\end{itemize}

\section*{Chapter 7: The Inhabited Library}

\begin{itemize}
\setlength{\itemsep}{0.5em}
\item The epigraph, ``The world is everything that is the case,'' is Wittgenstein, \textit{Tractatus Logico-Philosophicus} (1921), proposition 1, not 1.1 as it is sometimes cited; the wording is C.\,K.\ Ogden's 1922 translation.
\end{itemize}

\section*{Chapter 8: The Observer}

\begin{itemize}
\setlength{\itemsep}{0.5em}
\item The epigraph is Bryce DeWitt, ``Quantum Mechanics and Reality,'' \textit{Physics Today} 23 (1970); the branching reading traces to Hugh Everett III's original paper (\textit{Reviews of Modern Physics} 29, 1957).
\item ``The Self-Sampling Assumption (SSA)'' and ``the Self-Indication Assumption (SIA)'' were named and systematized by Nick Bostrom, \textit{Anthropic Bias} (2002). Appendix~\ref{app:anthropics} treats both, and the speed prior, in full.
\end{itemize}

\section*{Chapter 9: The Limits}

\begin{itemize}
\setlength{\itemsep}{0.5em}
\item The ``Boltzmann brain'' extends Ludwig Boltzmann's 1890s proposal that a low-entropy universe might itself be a rare fluctuation; its modern cosmological form is due to Albrecht and Sorbo (2004) and to Dyson, Kleban, and Susskind (2002).
\end{itemize}

\section*{Chapter 11: The Question of Continuation}

\begin{itemize}
\setlength{\itemsep}{0.5em}
\item ``This is quantum immortality'' names an implication most associated with Max Tegmark's ``quantum suicide'' thought experiment (\textit{Fortschritte der Physik} 46, 1998), with earlier versions from Hans Moravec (1987) and Bruno Marchal (1988). As the chapter says, it is a labeled implication, not a settled result.
\end{itemize}

\section*{Chapter 12: The Dissolution of Death}

\begin{itemize}
\setlength{\itemsep}{0.5em}
\item The epigraph, ``Death is nothing to us,'' is Epicurus, \textit{Letter to Menoeceus} (c.\ 300 BCE); it survives only because Diogenes Laertius quoted it whole in \textit{Lives of Eminent Philosophers}, Book X. No independent manuscript of Epicurus's writing on the point exists.
\end{itemize}

\section*{Chapter 13: Love Across Branches}

\begin{itemize}
\setlength{\itemsep}{0.5em}
\item The epigraph is Rilke, \textit{Letters to a Young Poet}. The ten letters were written 1902 to 1908, to a young cadet, Franz Xaver Kappus; 1929 is the book's posthumous publication date, three years after Rilke's death.
\end{itemize}

\section*{Chapter 14: The Pattern}

\begin{itemize}
\setlength{\itemsep}{0.5em}
\item The epigraph is David Hume, \textit{A Treatise of Human Nature} (1739 to 1740), Book I, Part IV, Section 6, ``Of Personal Identity.''
\item ``Derek Parfit argued that personal identity is not a deep further fact'' cites Parfit, \textit{Reasons and Persons} (1984), Part III. ``Thomas Metzinger argued that the self is a model the brain runs'' cites Metzinger, \textit{Being No One} (2003), restated for general readers in \textit{The Ego Tunnel} (2009).
\end{itemize}

\section*{Chapter 15: The Pattern Describing Itself}

\begin{itemize}
\setlength{\itemsep}{0.5em}
\item ``\textit{Il faut imaginer Sisyphe heureux}'' is the closing line of Albert Camus, \textit{Le Mythe de Sisyphe} (1942); Justin O'Brien's standard translation is \textit{The Myth of Sisyphus} (1955).
\item The epigraph is George Eliot, \textit{Middlemarch}, the Finale; 1872 is the completed four-volume edition, though the novel first appeared in installments beginning December 1871.
\end{itemize}
